\section{Introduction}
An embedded spiking controller must act before contact. A neural-rate criterion can accept a failed task, while altered activity can accompany success. Assessment requires a defined plant, causal actuation and a declared horizon.

Feedback couples faults to future sensory input. Equal seed labels need not align inputs when stopping or event consumption changes. Early termination also changes accumulated spikes: stopped, common-duration and fixed-horizon counts answer different resource questions. Pairing requires verified healthy capability in every direction.

We couple a fly-connectome SNN to synthetic pursuit avoidance. Assumed receptive fields and body coupling define an engineering stress model, not biological escape. Separate development and validation establish bilateral laws; new seed clusters cross law, direction and nested deletion dose at one second.

Our contribution is an executable paired protocol with source-bound registrations, shared input variates and masks, exact actuation and independent input/geometry replay. Prospective testing demonstrates feedback-dependent fault loss despite identical healthy success counts, and quantifies neural warning errors. Healthy reserves and actions differ: this is a whole-system comparison. Standard statistics and kinematics support bounded findings, not algorithmic novelty. Failed development and all test cells are retained.
